\newpage\EMPHASIZE{67. Conversion}

\textsc{Objectives}
\begin{itemize}
\item Use constructors for conversion
\item Overload conversion operators
\item Define constructors explicitly
\end{itemize}


\newpage\EMPHASIZE{Conversion using a constructor}

C++ will generate an automatic conversion using a constructor call if
necessary in order to find a best match for a method.

\begin{consolethree}[escapeinside=||]
class Int
{
public:
    Int(int a)
        : x_(a)
    {}
    void m(Int c) { std::cout << x_ << '\n'; }
private:
    int x_;
};
\end{consolethree}

\begin{consolethree}[escapeinside=||]
int main()
{
    Int c(0), d(1);
    c.m(d);
    |{\bf c.m(2);}|
    return 0;
}
\end{consolethree}

However, there is an exception: The constructor is \textbf{not} invoked
for the left-hand side of \textit{.} And \textit{->}:
\textit{3.m(1)} will not become \textit{C(3).m(1)}. You still might want to
write your own \textit{m(int)} for efficiency.

But there are problems with conversion using constructors:

\begin{itemize}
\item You cannot convert \textbf{to} built-in types (int, double, etc.)
  because they are not objects
\item You cannot convert from an object of class \textit{C} to an object of
  class \textit{D} (without changing the declaration of \textit{D})
\end{itemize}

So ...


\newpage\EMPHASIZE{Conversion operators}

Suppose \textit{D} is a type (class or built-in type like int, double,
...) and \textit{C} is a class. You can define a \textbf{conversion
operator} from \textit{C} to \textit{D}. This should look like

\begin{consolethree}[escapeinside=||]
class C
{
public:
    ...
    operator D();
    ...
};
\end{consolethree}

WARNING: Do not specify a return type. C++ knows that the return type
will be \textit{D}. Also, conversion operators are usually \textit{const}.

Let's make an addition to our \textit{Int} class

\begin{consolethree}[escapeinside=||]
class Int {
public:
    Int(int a)
        : x_(a)
    {}
    |{\bf operator int() const \{ return x\_; \}}|
private:
    int x_;
};
\end{consolethree}

\begin{consolethree}[escapeinside=||]
int main()
{
    Int c(0);
    std::cout << |{\bf (int) c}| << " " << |{\bf int(c)}|;
    return 0;
}
\end{consolethree}

\textit{(int) c} and \textit{int(c)} are the same as \textit{c.operator
int()}.

Here's another example:

\begin{consolethree}[escapeinside=||]
class D
{
public:
    D(int a, int b)
        : x_(a), y_(b)
    {}
    void print()
    {
        std::cout << "(" << x_ << "," << y_ << ")\n";
    }
private:
    int x_, y_;
};
\end{consolethree}

\begin{consolethree}[escapeinside=||]
class Int
{
public:
    Int(int a)
        : x_(a)
    {}
    operator D() { return D(x_,x_); }
private:
    int x_;
};
\end{consolethree}

\begin{consolethree}[escapeinside=||]
int main()
{
    Int c(2);
    ((D)c).print();
    D(c).print();
    return 0;
}
\end{consolethree}


\newpage\EMPHASIZE{Automatic Conversion}

C++ will automatically generate a call to conversion operators if
necessary to find a best match for a method.

\textbf{WARNING:} C++ will only generate one automatic conversion for
each value.

\textbf{WARNING:} Default conversions are used first (example
\textit{double()}).

C++ will use user-defined conversion only when necessary. No errors
during compilation.

\begin{consolethree}[escapeinside=||]
class D
{
public:
    D(int a, int b)
        : x_(a), y_(b)
    {}
private:
    int x_, y_;
};
\end{consolethree}

\begin{consolethree}[escapeinside=||]
class Int
{
public:
    Int(int a)
        : x_(a)
    {}
    operator D()
    {
        std::cout << "Int::D()\n";
        return D(x_,x_);
    }
    void m(D d)
    {
        std::cout << "m(D)\n";
    }
private:
    int x_;
};
\end{consolethree}

\begin{consolethree}[escapeinside=||]
int main()
{
    Int c0(2), c1(3);
    c0.m(c1);
    return 0;
}
\end{consolethree}

\begin{consolethree}[escapeinside=||]
class E {};

class D
{
public:
    operator E()
    {
        ...
    }
};

class C
{
public:
    operator D()
    {
        ...
    }
};

int main()
{
    E e1 = D(C());
    E e2 = D();
    |{\bf E e3 = C();}|
    return 0;
}
\end{consolethree}

\textbf{ERROR:} \textbf{Two} automatic conversions needed, \textit{C} to
\textit{D} to \textit{E}. i.e. C++ will \textbf{not} convert

\textit{E e3 = C()}

to

\textit{E e3 = E(D(C()))}

Here's another example:

\begin{consolethree}[escapeinside=||]
class D {};

class C
{
public:
    operator D() { std::cout << "C::D()\n"; }
};

void f(D d) { std::cout << "f(D)\n"; }

void f(double d) { std::cout << "f(double)\n"; }

int main()
{
    f(1);
    return 0;
}
\end{consolethree}

Notice that the default conversion to a \textit{double()} was used instead
of the \textit{C::D()} conversion.


\newpage\EMPHASIZE{Advice}

Do not have too many conversion operators. Doing so might lead to
ambiguities. For instance, suppose you have conversion operators from
\textit{C} to \textit{D} and vice versa; you also have \textit{operator+} in
\textit{C} and \textit{D}. If you run the following code

\begin{consolethree}[escapeinside=||]
C c;
D d;
c + d;
\end{consolethree}

you'll see that your compiler does not know if it should do

\textit{D(c) + d} using \textit{D::operator+}

or

\textit{c + C(d)} using \textit{C::operator+}

Here's another tip: If you have two classes \textit{C} and
\textit{D} where \textit{C} is ``included'' in D, then it's
better to have automatic conversion from \textit{C} to \textit{D}. Consider
the following example:

Let's say you have an \textit{Int} class and a
\textit{Fraction} class. You should have a \textit{operator Fraction()}
declared in the \textit{Int} class. Note that you can achieve the same
thing by having a constructor of the form:

\textit{Fraction(const Int \&) const;}


\newpage\EMPHASIZE{Explicit}

You can prevent automatic type conversion via constructors.
Here's an example:

\textit{\textbf{explicit} Fraction(int);}

In this case, C++ will \textbf{not} perform automatic type conversion
from an \textit{int} to a \textit{Fraction}.

Here's an example of what you've already seen before:

\begin{consolethree}[escapeinside=||]
class C
{
public:
    C(int i)
    {}
};

void f(C c)
{
}

int main()
{
    f(1);
    return 0;
}
\end{consolethree}

You should already know that this will not produce any errors;
\textit{f(1)} automatically becomes \textit{f(C(1))}. However, if you run
this

\begin{consolethree}[escapeinside=||]
class C
{
public:
    |{\bf explicit}| C(int i)
    {}
};

void f(C c)
{
}

int main()
{
    f(1);
    return 0;
}
\end{consolethree}

You'll get an error! The compiler does not know what to
do because we did not define a way for an \textit{int} to be converted to
a \textit{C} object.

Note that

\begin{itemize}
\item Only constructors can be made explicit.
\item You cannot make type conversion operators which are not a constructor
  explicit.
\end{itemize}

Here's an example of the second restriction:

\begin{consolethree}[escapeinside=||]
class C
{
    explicit C(int); // OK
    |{\bf explicit}| operator int() const; // |{\bf BAD!}|
};
\end{consolethree}

\textbf{In the future}, C++ will allow all type conversion to be made
\textit{explicit}. (NOTE: C++11 (i.e., C++ 2011) supports \textit{explicit}
for type conversion.)

Now, suppose the \textit{Fraction} class has the constructor
\textit{Fraction(int)} and \textit{operator int()}. If we run

\textit{Fraction r(1, 2);}

\textit{std::cout << r + 1 << '\textbackslash n';}

Which is used? \textit{int(r) + 1} or \textit{r + Fraction(1)}? \textbf{If} \textit{operator int()}
is explicit. Then the above would execute

\textit{r + Fraction(1)}

\end{document}
